\chapter{Пробная реализация единой сущности}

Для проверки понятийного каркаса вводится пробная реализация
\begin{equation}
 \mathfrak S_{\mathrm{trial}}=
 \bigl(\mathcal B,\Phi,\mathcal L_{\mathrm{int}},\mathcal C\bigr),
\end{equation}
где $\mathcal B$ обозначает опорный каркас, $\Phi$ --- глобальный фазовый
сектор, $\mathcal L_{\mathrm{int}}$ --- внутреннюю нагрузку, а $\mathcal C$ ---
правила их согласования.

Эта запись не является окончательным определением $\mathfrak S$. Она нужна,
чтобы проверить, можно ли одним объектом организовать разные типы наблюдаемых
без независимого правила для каждого из них.

\section{Опорный каркас}

Каркас хранит данные, устойчивые при слабых фазовых и внутренних изменениях:
топологический класс, ведущую спектральную плотность, размерность носителя и
глобальные условия допустимости.

\section{Фазовый сектор}

Фазовый сектор различает ветви, которые нельзя устранить локальной сменой
координат. Его кандидатами являются голономии, спиновые структуры и
ориентированные глобальные фазы.

\section{Внутренняя нагрузка}

Внутренний сектор отвечает за кратности, тонкие расщепления и операторные
связи, не определяемые одним каркасом. Его данные не должны вводиться как
произвольный список масс или коэффициентов.

\section{Правила согласования}

Объект $\mathcal C$ задаёт, какие сочетания секторов допустимы. Именно здесь
должны возникать смешанные наблюдаемые. Если $\mathcal C$ заменяется набором
независимых поправок, пробная реализация не подтверждает единый источник.

\chapter{Локальный пример чтения}

Рассмотрим семейство операторов
\begin{equation}
 D(\alpha,a)=D_{\mathcal B}+A_{\Phi}(\alpha)+Y_{\mathrm{int}}(a)
 +M_{\times}(\alpha,a).
\end{equation}

Здесь $\alpha$ отмечает глобальную фазовую ветвь, а $a$ --- внутренний канал.
Последний член отвечает за несводимое смешение.

\section{Чистые пределы}

При $Y_{\mathrm{int}}=M_{\times}=0$ сравниваются только фазовые ветви одного
каркаса. При $A_{\Phi}=M_{\times}=0$ исследуется внутренняя структура без
глобального фазового расщепления.

Эти пределы служат контрольными исключениями. Они показывают, какие признаки
исчезают при удалении каждого сектора.

\section{Смешанный остаток}

Член $M_{\times}$ считается содержательным, если его нельзя поглотить
переопределением $A_{\Phi}$ или $Y_{\mathrm{int}}$ и если он сохраняется при
допустимой смене базиса.

\section{Наблюдаемое чтение}

Для спектрального функционала $F$ пробное наблюдаемое записывается как
\begin{equation}
 O(\alpha,a)=\operatorname{Tr}F\bigl(D(\alpha,a)\bigr).
\end{equation}

Секторная атрибуция определяется изменением $O$ при последовательном
исключении $A_{\Phi}$, $Y_{\mathrm{int}}$ и $M_{\times}$.

\chapter{Полуколичественная реконструкция}

До получения полного оператора допускается ограниченная схема чувствительности.
Пусть
\begin{equation}
 r_a(O)=\left|\frac{\partial O}{\partial\lambda_a}\right|_{\lambda=0}
\end{equation}
измеряет отклик на заранее нормированную деформацию сектора $a$.

\section{Требование общей нормировки}

Все $r_a$ должны вычисляться в одной метрике пространства деформаций. Нельзя
независимо масштабировать направления так, чтобы получить желаемую секторную
картину.

\section{Нормированные веса}

Вводятся
\begin{equation}
 w_a(O)=\frac{r_a(O)}{\sum_b r_b(O)+\epsilon},
 \qquad \sum_a w_a\simeq1.
\end{equation}

Регуляризатор $\epsilon$ фиксируется до анализа и должен исчезать в допустимом
пределе. Веса являются характеристиками выбранной схемы чувствительности, а
не новыми фундаментальными постоянными.

\section{Область применимости}

Схема пригодна для сравнения режимов и проверки устойчивости, но не для
абсолютного предсказания масс, энергий или вероятностей. Размерные выводы
требуют полного действия и физической нормировки.

\chapter{Инвариант некаркасности}

Первой координатой карты служит отношение некаркасного и каркасного откликов:
\begin{equation}
 \Xi(O)=
 \frac{r_{\mathrm{phase}}(O)+r_{\mathrm{int}}(O)+r_{\mathrm{mix}}(O)}
 {r_{\mathrm{car}}(O)+\epsilon}.
\end{equation}

\section{Толкование}

При $\Xi\ll1$ наблюдаемое устойчиво каркасно. Значение порядка единицы
указывает на переходную область. Большое $\Xi$ означает, что одной геометрии
или топологии носителя недостаточно.

\section{Проверка устойчивости}

Классификация по $\Xi$ принимается только при сохранении порядка величины под
допустимыми деформациями и при малой зависимости от $\epsilon$.

\section{Ограничение}

$\Xi$ не является физической наблюдаемой до вывода канонической метрики
деформаций. Она служит сравнительным показателем внутри одного протокола.

\chapter{Инвариант внутренней ориентации}

Вторая координата различает фазовый и внутренний вклады:
\begin{equation}
 \Upsilon(O)=
 \frac{r_{\mathrm{phase}}(O)-r_{\mathrm{int}}(O)}
 {r_{\mathrm{phase}}(O)+r_{\mathrm{int}}(O)+\epsilon}.
\end{equation}

\section{Области значений}

Значение $\Upsilon\approx1$ соответствует фазовой ориентации,
$\Upsilon\approx-1$ --- внутренней, а $\Upsilon\approx0$ --- отсутствию
устойчивого преобладания между ними.

\section{Смешанный сектор}

Нулевое значение $\Upsilon$ само по себе не доказывает смешанный механизм:
возможна простая компенсация двух независимых вкладов. Поэтому дополнительно
проверяется ненулевой и несводимый $r_{\mathrm{mix}}$.

\section{Совместное чтение}

Пара
\begin{equation}
 \mathfrak R(O)=\bigl(\Xi(O),\Upsilon(O)\bigr)
\end{equation}
различает степень ухода от каркаса и направление этого ухода. Она полезна для
сравнения классов наблюдаемых, но сохраняет методический статус до строгого
вывода нормировки.

Следующий и последний блок должен содержать иллюстративные схемы, заключение,
послесловие после проверки вторым томом, перспективу и окончательную
заморозку статуса.